\section{Background and Motivation}
\label{sec:motivation}

\subsection{Primer on Multi-Access Edge Computing}
\label{sec:primer}

Multi-Access Edge Computing (MEC) brings computational resources closer to end users by deploying compute infrastructure close to cellular base stations~\cite{etsi-mec}.
The MEC architecture consists of three key components: the RAN that manages wireless spectrum and radio resource scheduling, edge servers that provide CPU and GPU resources, and the core network that connects the RAN and the servers.
In the typical MEC request-response model, LC applications offload compute-intensive tasks from client devices to edge servers.
Clients send requests via RAN uplink, edge servers process them locally, and responses return via RAN downlink.

\begin{table}[]
\setlength{\tabcolsep}{4pt}
\renewcommand{\arraystretch}{1.0}
\footnotesize
\centering
\begin{tabular}{@{}p{2.1cm}p{2.0cm}p{0.8cm}p{1.2cm}p{1.2cm}@{}}
\toprule
\textbf{Applications} & \textbf{Offloaded Task} & \textbf{SLO} & \textbf{UL/DL Load} & \textbf{Compute Resource} \\ \midrule
Smart stadium~\cite{www-orange-5gstadium} & Video transcoding & 100ms & High/High & CPU \\
Augmented reality~\cite{mao2019delay, padmanabhan2023gemel} & Object detection & 100ms & Med/Low & GPU \\
Video conferencing~\cite{www-zoom-latency, www-wikipedia-audio-latency} & Super resolution & 150ms & Low/High & GPU \\ \bottomrule
\end{tabular}
\caption{Examples of MEC applications evaluated in this paper, each with distinct SLO, network (uplink/downlink) load, and compute resource requirements. Details are described in~\autoref{sec:eval-setup}.}
\label{tbl:apps}
\vspace{-1em}
\end{table}

\autoref{tbl:apps} shows the MEC applications we focus on in this paper, each with distinct SLO, network load, and compute resource requirements.
We target realistic MEC deployments where multiple user equipment (UE) devices (\eg smartphones, cameras, AR headsets) simultaneously offload compute-intensive tasks to edge servers.

However, as we demonstrate next, the reality of current MEC deployments falls short of these performance promises.

\subsection{Unpredictable Performance of MEC}
\label{sec:unpredictable}

We benchmarked MEC deployments across three cities in the US, South Korea, and China (Dallas, Seoul, and Nanjing) to quantify this performance gap.

\mypara{Measurement setup.}
We deployed the smart stadium and augmented reality applications from~\autoref{tbl:apps} using a laptop client connected to a 5G smartphone hotspot and edge VMs provisioned through edge service providers (\eg AWS in the US).
Each application runs in isolation (\ie no contention on the VM), generates 10,000 requests, and we measure end-to-end latency from request transmission to complete response reception.
For brevity, we focus our analysis on the smart stadium application, presenting results from all three cities that employ different combinations of cellular operators and cloud providers.
We observe similar trends for the augmented reality application across all cities (\autoref{appnd:ar}).

\begin{figure}[]
    \centering
    \includegraphics[width=\columnwidth]{figures/measurement-contention-results/ran-contention/video-transcoding/e2e_latency_cdf.pdf}
    \tightcaption{End-to-end latency for the smart stadium application without edge resource contention across MEC deployments in three cities. The dotted red line indicates the SLO.}
    \label{fig:motive-e2e-transcoding-ran}
\end{figure}

\mypara{Results.}
\autoref{fig:motive-e2e-transcoding-ran} reveals unpredictable performance characterized by high tail latencies across all deployments, even without compute resource contention on the edge server.
Specifically, 7\%, 20\% and 47\% of requests exceed their SLO requirements in \dallas, \nanjing, and \seoul, respectively, during low network activity periods (measured at 2am).
Although median latencies remain below the SLO threshold, the P95 and P99 latencies are substantially higher, resulting in an inconsistent user experience.
This problem intensifies under higher network load: when additional UE devices access the 5G network and increase contention for RAN resources (\dallas-Busy), even the median latency exceeds the SLO requirements.


\subsection{Root Causes of Unpredictability}
\label{sec:root-cause}

To understand the sources of this unpredictability, we investigate what contributes to the tail latency in both the RAN and the edge server.

\subsubsection{RAN Resource Contention}
\label{sec:wireless-bottlenecks}

To analyze RAN-induced latency variability, we decompose end-to-end latency into uplink and downlink transmission phases.
Using Precision Time Protocol (PTP)~\cite{ptp-std} for clock synchronization between client device and edge server via a stable out-of-band wired connection, we developed a synthetic application that  measures uplink and downlink latency separately while varying request and response sizes.
Specifically, we measured the time to receive complete requests at the server (uplink) and complete responses at the client (downlink) while varying data size. 

\begin{figure}[]
    \centering
    \includegraphics[width=\columnwidth]{figures/measurement-decompose-results/total_latency_boxplot.pdf}
    \tightcaption{Network latency variability for uplink and downlink transmissions across different data sizes in \dallas.
 Data size refers to request size for uplink and response size for downlink.}
    \label{fig:motive-decompose}
\end{figure}

\autoref{fig:motive-decompose} reveals an asymmetry that explains MEC's unpredictable performance: 
uplink latency exhibits high variability, especially for larger request sizes, while downlink latency remains stable.
This reflects how cellular networks provision fewer uplink slots than downlink slots, causing higher contention for uplink transmissions.
We observe similar trends across all measured cities (\autoref{appnd:latency-variability}).

To understand the specific mechanisms causing uplink jitter, we leveraged our srsRAN 5G~\cite{srsran}-based testbed for in-depth analysis (details in~\autoref{sec:eval-setup}), since we have no visibility into the internals of public cellular networks.
The srsRAN stack employs the proportional fair (PF) scheduling algorithm~\cite{kelly1997charging,jalali2000data} used in commercial deployments, allowing us to emulate real-world scheduling behavior.

We ran the smart stadium application on the testbed and monitored MAC layer state by collecting buffer status reports (BSRs) sent from each UE.
As a background workload, we also deployed five file transfer UEs.
BSRs indicate remaining data in UE-side transmission buffers, providing direct insight into whether the scheduler allocates sufficient RAN resources\footnote{We use the term RAN resources to refer to uplink and downlink spectrum resources managed by the MAC scheduler (\ie Physical Resource Blocks in 5G).} to meet application demands.

\begin{figure}[]
    \centering
    \includegraphics[width=\columnwidth]{figures/transcoding-bsr/request_bsr.pdf}
    \tightcaption{
        Smart stadium UE's uplink buffer status changes over time. 
        300 KB is the maximum for BSR from UE to the RAN, which means UE may buffer more than 300 KB.
    }
    \label{fig:uplink-bsr}
\end{figure}

\autoref{fig:uplink-bsr} shows persistent non-zero BSR (>1s), indicating uplink starvation from PF scheduling.
This resource starvation directly translates to the observed uplink jitter; inadequate resource allocation causes unpredictable buffering delays that lead to SLO violations.

The root cause is PF schedulers' lack of SLO awareness.
They prioritize UEs with relatively better channel conditions to their historical average throughput, balancing fairness and efficiency without considering SLO requirements.
When multiple UEs compete for uplink resources, LC applications cannot obtain timely allocations, missing SLO requirements even when aggregate bandwidth is sufficient.

\begin{figure}[]
    \centering
    \includegraphics[width=\columnwidth]{figures/measurement-contention-results/compute-contention/video-transcoding/Dallas/e2e_latency_cdf.pdf}
    \tightcaption{End-to-end latency for smart stadium under different levels of compute resource contention in \dallas. The dotted red line indicates the SLO.}
    \label{fig:motive-e2e-transcoding-compute}
\end{figure}

\subsubsection{Edge Compute Resource Contention}
\label{sec:compute-bottlenecks}

While wireless resource allocation causes unpredictable uplink performance, edge compute resources present an additional bottleneck in MEC deployments.
To quantify the impact of compute resource contention, we emulate competing offloaded tasks by running a CPU stressor (stress-ng~\cite{stress-ng}) with varying CPU usage levels.
\autoref{fig:motive-e2e-transcoding-compute} shows that compute resource contention significantly contributes to unpredictable performance of the smart stadium application.
As CPU load increases, tail latencies grow substantially, confirming that inadequate management of edge compute resources causes processing delays that cascade into SLO violations.
This problem extends to GPU-intensive applications and other MEC deployments (\autoref{appnd:compute-contention}).


\subsection{Prior Approaches and Limitations}
\label{sec:prior}

Having identified the root causes of unpredictable MEC performance, we now examine existing resource management approaches and discuss why they fall short for our setting.

\mypara{MEC-specific resource management.}
Recent MEC resource management proposals suffer from impractical deployment requirements and limited scope as we show in~\autoref{sec:eval}.
\tutti{}~\cite{xu2022tutti} assumes explicit coordination between the RAN  and edge applications, requiring servers to notify the RAN upon receiving requests and transport protocol modifications to identify ``first packets.''
This is impractical because RAN infrastructure and edge servers are typically managed by different entities (\eg Verizon operates the RAN while AWS provides edge compute).
Even setting deployment issues aside, \tutti{} fails to timely accelerate LC requests due to feedback delays between edge servers and RAN.
Moreover, \tutti{} only supports homogeneous applications with identical SLOs.
\arma{}~\cite{arma:mobisys25} similarly requires coordination between edge servers and RAN schedulers, is narrowly tailored to video analytics, and allows non-LC applications to block LC ones when their uplink bandwidth usage is high.
Both systems focus solely on RAN scheduling while ignoring edge compute resource management for co-existing heterogeneous applications.

\mypara{SLO-aware resource management in clouds.}
A natural approach for managing compute resources is to adapt existing SLO-aware resource management solutions from cloud environments.
However, these solutions are ill-suited for our target LC applications.
PARTIES~\cite{parties:asplos19} reactively adjusts server resource partitions based on SLO feedback from clients, but fails in MEC where wireless feedback delays mean multiple requests miss deadlines before adjustments take effect.
Caladan~\cite{caladan:osdi20} employs proactive CPU core allocation without requiring client feedback but requires extensive application modifications.
ML-based schedulers like OSML~\cite{liu2023intelligent} incur inference overhead and operate at second-level intervals, making them too slow to react within the tens of milliseconds required for LC applications.
Similarly, UFO~\cite{peng2024ufo} relies on changes in system-wide metrics (\eg CPU utilization, scheduling frequency) captured at second-level intervals to infer potential SLO violations and makes scheduling decisions, which reacts too slowly for LC applications. 
Vessel~\cite{lin2024fast} assumes all processes share one memory space and relies on hardware-assisted core allocation, limiting its applicability in multi-tenant MEC settings and beyond CPU scheduling.

\mypara{End-to-end rate control mechanisms.}
End-to-end congestion control mechanisms~\cite{carlucci2017congestion,zhu2020network,johansson2017self} are ill-suited for LC applications in MEC environments.
Congestion control requires hundreds of milliseconds to converge, causing SLO violations before it can react to uplink congestion.
Also, the 5G uplink channel quality fluctuates rapidly due to limited UE transmission power and varying user counts, making it nearly impossible for congestion control to stabilize.
\modified{Even approaches that leverage 5G control-channel information (\eg PBE-CC~\cite{xie2020pbe}) to react quickly run solely at the end host and cannot directly influence RAN uplink scheduling.
Under severe wireless contention, the RAN may allocate minimal or no uplink resources to a UE, so sender-side adaptation alone is insufficient to satisfy SLO requirements.}
Additionally, many LC applications rely on constant or variable bitrate encoding to maintain video quality~\cite{Fortinet2020_IPSurveillanceCameraBandwidth,Sony_BRCX400_Specs}, where reactive rate adaptation directly degrades quality.

\modified{\mypara{URLLC.}
5G Ultra-Reliable Low-Latency Communication (URLLC), a service category designed for mission-critical applications, achieves low latency and high reliability through conservative radio configurations and dedicated resource reservations~\cite{esti-5g-arch}.
However, this approach reduces spectrum efficiency: reserved resources often remain underutilized, and reliability-enhancing techniques consume additional radio resources without commensurate benefit.
The inefficiency is particularly pronounced for uplink-heavy workloads where resource demand is dynamic, making static reservations ill-suited for our setting.}